\section{Universal Data from the Dynamics}
\label{sec:cft}

The premise of this thesis is that the real-time dynamics of a lattice model, however microscopic and messy, carries sharp universal information when the model is critical, and that this information can be read directly off a simulation. The bridge is conformal field theory. At a continuous critical point a lattice model flows to a \ac{CFT}, a continuum theory so constrained by its symmetry that a handful of its predictions are fixed by a single number, the central charge $c$, together with a discrete list of operator dimensions. These numbers are universal: every model in the same universality class must reproduce them, which is exactly what makes them good targets for a numerical method. This section explains how a dynamical quantity, the Loschmidt echo, is mapped onto such a theory; what universal data the theory then predicts; and, crucially, how that universal data is separated from the non-universal lattice constants that ride along with it in any real simulation. The concrete predictions collected at the end are what Section~\ref{sec:results} measures.

We keep the conformal machinery to the minimum the method needs. The few derivations we use are collected in Appendix~\ref{app:cft}, and standard material is referred to~\cite{difrancesco1997, bpz1984, cardy1989}.

\subsection{From the Loschmidt echo to a boundary CFT}
\label{sec:cft:boundary}

The object we compute is the Loschmidt echo introduced in Section~\ref{sec:intro},
\begin{equation}
    \mathcal{L}(T) = \bra{\Psi_0}\, e^{-iHT} \ket{\Psi_0},
    \label{eq:echo}
\end{equation}
the amplitude for the system, prepared in the product state $\ket{\Psi_0}=\ket{X+}^{\otimes N}$, to return to it after evolving for a time $T$ with the critical Hamiltonian. To bring conformal field theory to bear, we pass to imaginary time, $t\to-i\tau$: the echo becomes a Euclidean path integral over a slab whose two ends are both fixed to $\ket{\Psi_0}$. In \ac{CFT} language, a state that can terminate a path integral is a conformal boundary state, and the admissible ones, the Cardy states~\cite{cardy1989}, are the allowed boundary conditions of the theory. The picture that results is the one everything below rests on: the echo is the partition function of the critical theory on a long strip whose two edges carry the boundary condition set by $\ket{\Psi_0}$, and whose width is the (imaginary) evolution time.

Two features of this strip recur throughout the numerics. First, a sharp product state has structure down to the lattice spacing and is not, by itself, a smooth conformal boundary state; we therefore pre-cool it in imaginary time for a short interval $\beta_0$, which on the lattice amounts to $N_\beta = 2\beta_0/\delta t$ imaginary-time sites split between the two ends. The strip width becomes $T+2\beta_0$, and the conformal predictions are recovered as $\beta_0$ is made small. This UV regulator is not optional: without it the boundary condition is dirty and the universal signal is spoiled. Second, which boundary condition we obtain depends on $\ket{\Psi_0}$. Each initial state realizes a particular Cardy state, whose lightest boundary operator has a dimension $x_1$ that controls how fast finite-size corrections die away. For the Ising class the disordered (free) boundary has $x_1=\tfrac12$ and converges quickly, whereas the ordered (fixed) boundary has $x_1=2$ and converges slowly; in the $\sigma^x$-field convention of our model the polarized $\ket{X+}$ realizes the free boundary, which is why we quench from it.

The step that turns this into a method is to read the strip sideways. In the Euclidean theory the two coordinates enter on the same footing---nothing intrinsically distinguishes ``space'' from ``imaginary time''---so we are free to quantize along space and propagate along time, or the other way round. We take the other way round: we build a transfer matrix $\mathcal{E}$ that advances the contraction along the temporal direction, acting on states that live on a time-like slice. The eigenproblem of that single matrix, as the next subsection shows, contains the entire universal content, and the tensor-network machinery of Section~\ref{sec:transverse} is what evaluates it.

\subsection{What the continuum theory gives us}
\label{sec:cft:content}

At the critical point the correlation length diverges, the model loses every intrinsic scale, and its long-distance physics is described by a scale-invariant continuum theory. In two dimensions, locality promotes scale invariance to invariance under all angle-preserving maps, the conformal transformations, which constrains the theory almost completely (Appendix~\ref{app:cft:maps};~\cite{bpz1984, difrancesco1997}). For our purposes this abstract statement has three concrete payoffs, each a channel through which a universal number reaches the transfer matrix.

The first is the central charge $c$ itself, the single label of the universality class ($c=\tfrac12$ for Ising). It is a conformal anomaly: confining the theory to a strip of transverse size $\beta$ costs a universal ground-state energy proportional to $c$. Quantizing on that strip, the generator of translations along it is~\cite{carignano_tagliacozzo2025}
\begin{equation}
    \hat H = \frac{\pi}{\beta}\Big(L_0 - \frac{c}{24}\Big),
    \label{eq:cft:casimir}
\end{equation}
where $L_0$ is the scaling generator. Its two terms carry the two kinds of universal data at once: the $-c/24$ term is the Casimir anomaly, so isolating it measures $c$; the $L_0$ term stacks the operator dimensions on top of that offset, as the excited levels.

The second payoff is the operator content. A \ac{CFT} is organized around primary operators, each with a scaling dimension
\begin{equation}
    x = h + \bar h
    \label{eq:cft:dim}
\end{equation}
that fixes the power law of its correlations; the list of dimensions is a rigid fingerprint of the theory. Because our transfer matrix is essentially $\mathcal{E}=e^{-\delta t\,\hat H}$, Equation~\eqref{eq:cft:casimir} hands us its spectrum directly~\cite{carignano_tagliacozzo2025},
\begin{equation}
    \mathcal{E} = \exp\!\Big[-\big(-\kappa + \pi L_0\big)\cdot(\text{geometry})\Big],
    \qquad \kappa = -\frac{\pi c\,\delta t}{24},
    \label{eq:cft:transfer}
\end{equation}
a dominant eigenvalue whose modulus encodes the Casimir term (hence $c$), and a tower of subleading eigenvalues whose gaps to it encode the dimensions $x$ (Appendix~\ref{app:cft:spectrum}). Diagonalizing one matrix therefore reads off both the central charge and the operator content.

The third payoff is entanglement. Computing a R\'enyi entropy in a \ac{CFT} reduces to a correlation function of twist operators, themselves primaries of dimension~\cite{calabrese_cardy2004}
\begin{equation}
    \Delta_n = \frac{c}{12}\Big(n - \frac{1}{n}\Big),
    \label{eq:cft:twist}
\end{equation}
so the entropy is a logarithm whose slope is proportional to $c$ (Appendix~\ref{app:cft:twist}). This is the mechanism behind every entropy-based measurement in this thesis, in equilibrium (Section~\ref{sec:annni}) and in time. The R\'enyi order $n$ enters through a known prefactor; for a boundary (half-line) cut, which is the situation of our temporal entropies, the coefficient of the logarithm is
\begin{equation}
    \frac{c}{12}\,\frac{n+1}{n},
    \qquad \text{i.e.}\quad \frac{c}{6}\ (n\to1), \qquad \frac{c}{8}\ (n=2),
    \label{eq:cft:coeff}
\end{equation}
half the familiar bulk value. We fix which member of this family applies when we define the entropy our simulations compute, in Section~\ref{sec:transverse}.

\subsection{Universal versus non-universal constants}
\label{sec:cft:universal}

Everything in the previous subsection is exact in the continuum. A simulation, however, never sees the pure continuum: it works at a finite evolution time $T$, a finite bond dimension $\chi$, a finite Trotter step $\delta t$, and a nonzero regulator $\beta_0$, on a lattice that carries its own microscopic constants. The measured quantities are therefore mixtures of the universal numbers we want and non-universal constants we do not. Separating the two is the methodological core of this work, and it is worth setting out before the predictions, because it dictates the form in which each prediction is written.

The universal quantities, fixed by the universality class and nothing else, are the central charge $c$, the scaling dimensions $x_i$, the logarithmic coefficients $c/6$ and $c/8$, the Casimir coefficient $\kappa=-\pi c\,\delta t/24$, and the pattern by which the leading eigenvalue moduli become degenerate. The non-universal constants are everything the lattice and the algorithm bring in: the regulator $\beta_0$; the overall per-column normalization of the transfer matrix, which sets the modulus of its dominant eigenvalue to a model-dependent value $|\lambda_0|\approx1.5$ rather than one; the additive offset $s_0$ of the entropy; the Trotter step $\delta t$; and the sound velocity $v$. The velocity deserves a separate word. It is genuinely non-universal---it depends on $p$ and grows with frustration (Section~\ref{sec:annni:velocity})---yet it is measurable independently from the equilibrium dispersion, and it is the bridge that turns a universal phase into a universal dimension, since every eigenvalue phase appears in the combination $\pi x/(vT)$. A universal dimension is only accessible once the velocity is in hand.

Four devices isolate the universal part, and each shapes one of the predictions that follow.

The first is the ratio of gaps. A single eigenvalue gap carries a universal dimension but also the non-universal $\beta_0$ and normalization; a ratio of two gaps cancels them, so that $(x_i-x_0)/(x_j-x_0)$ is pure operator content, independent of $\beta_0$, of the normalization, and even of the velocity. This is the cleanest universal read available, and the one we trust most for the operator content.

The second is calibration against the integrable point. The entropy slope measures $c$, but at finite $(T,\chi,\delta t)$ it is not $c/6$ exactly: it is $c/6$ times a prefactor that collects the finite-time conformal correction, the truncation, and the Trotter error. Writing the leading behaviour as
\begin{equation}
    \mathrm{slope}(p; T, \chi, \delta t) = \frac{c(p)}{6}\,K(T,\chi,\delta t,\beta_0) + \text{(subleading)},
    \label{eq:cft:calib}
\end{equation}
the prefactor $K$ is the same function for every model in the class: the correction it collects is a property of the method and of the finite-time conformal correction, whose form is universal, not of the particular coupling. Since the integrable point $p=0$ has $c=\tfrac12$ exactly and is run through the identical pipeline at matched $(T,\chi,\delta t,\beta_0)$, taking the ratio cancels $K$,
\begin{equation}
    c(p) = \frac{1}{2}\,\frac{\mathrm{slope}(p)}{\mathrm{slope}(0)},
    \label{eq:cft:calibratio}
\end{equation}
and returns the universal $c(p)$. It is worth being explicit that this is not a low-$T$ coincidence, because the distinction is what makes the technique trustworthy. A coincidence would appear at one value of $T$ and disappear at others; the calibrated $c(p)$, by contrast, is flat across the whole clean window (Section~\ref{sec:results}). The raw slope drifts with $T$ precisely because $K(T)$ drifts, but $K(T)$ drifts the same way for $p=0$ and for $p\neq0$, so their ratio stands still. The only thing that does not cancel is the difference between the two couplings in the subleading amplitude of Equation~\eqref{eq:cft:calib}, which is small in the clean window and is exactly what sets the error bar on $c(p)$. The calibration removes a shared, class-level systematic exactly, and leaves a small, bounded residual: a controlled cancellation, not luck.

The third device is to read phases rather than moduli. In Equation~\eqref{eq:cft:transfer} the universal data sits in the phases of the eigenvalues, while the modulus additionally carries the non-universal normalization. Reading the phases is therefore reading universal information directly, and the modulus enters only through differences and ratios. This is also why the signature of emergent dual unitarity below is the constancy of the eigenvalue modulus, not its value: the value is non-universal, its becoming constant is not.

The fourth is to treat the finite-time corrections as signal. The corrections lumped into $K$ are not noise; they are themselves universal, with a form fixed by the \ac{CFT}~\cite{boucomas2026}. One can therefore either cancel them, through a ratio or the calibration above, or model them explicitly and read $c$ absolutely; both routes appear in Section~\ref{sec:results}. A related subtlety, a constant branch offset that must be pinned before the small universal Casimir term becomes visible, is spelled out in Appendix~\ref{app:eq3}.

\subsection{The universal predictions}
\label{sec:cft:temporal}
\label{sec:emergent}

Reading the strip of Section~\ref{sec:cft:boundary} sideways yields, for the generalized entropy along a temporal cut, the closed form~\cite{carignano_tagliacozzo2025}
\begin{equation}
    S^{\mathrm{gen}}(t) = s_0 + \frac{i\pi c}{12}
      + \frac{c}{6}\,\log\!\Big[\frac{2T}{\pi}\sin\frac{\pi t}{T}\Big],
    \label{eq:cft:sgen}
\end{equation}
whose decisive feature is that, the separation being time-like, the entropy grows only logarithmically in $T$ rather than linearly. That logarithmic growth is what makes the temporal contraction efficient in the first place. Together with the transfer spectrum of Equation~\eqref{eq:cft:transfer}, this gives four universal predictions, which Section~\ref{sec:results} confronts with the simulations. Each is written in the form its universal content dictates (Section~\ref{sec:cft:universal}); throughout, $\lambda_i=\log(-\tau_i)$ denotes the log of the transfer eigenvalue $\tau_i$, so phase differences are read as $\arg(-\tau)$.

\begin{enumerate}
    \item The generalized-entropy profile. Splitting Equation~\eqref{eq:cft:sgen} into real and imaginary parts,
    \begin{equation}
        \mathrm{Re}\,S^{\mathrm{gen}}(t) = s_0 + \frac{c}{6}\log\!\Big[\frac{2T}{\pi}\sin\frac{\pi t}{T}\Big],
        \qquad
        \mathrm{Im}\,S^{\mathrm{gen}} = \frac{\pi c}{12}.
        \label{eq:cft:chord}
    \end{equation}
    The real part is a symmetric dome; plotted against the chord variable $\log[\tfrac{2T}{\pi}\sin\tfrac{\pi t}{T}]$ it is a straight line whose slope is $c$ times the coefficient of Equation~\eqref{eq:cft:coeff}. Since that slope carries the non-universal prefactor $K$, we read $c$ from it by calibration (Equation~\eqref{eq:cft:calibratio}), or absolutely once the finite-time correction is restored. The imaginary part is a flat plateau: $\pi c/12$ for the $n\to1$ entropy, and $\tfrac{c}{12}(1+\tfrac1n)\tfrac{\pi}{2}$, that is $\pi c/16$ at $n=2$, for the R\'enyi generalization~\cite{boucomas2026}---a second read of $c$ from the same data.

    \item The central charge from the leading eigenvalue. The imaginary part of the dominant eigenvalue obeys
    \begin{equation}
        \frac{\mathrm{Im}\,\lambda_0}{T} = a_0 + \frac{\kappa}{T^2} + \frac{a_4}{T^4} + \cdots,
        \label{eq:cft:eq3}
    \end{equation}
    so the universal Casimir coefficient $\kappa$ sits in the $1/T^2$ term, underneath the non-universal extensive term $a_0$. Isolating it gives $c$ through the phases alone, independently of the entropy. It requires a nonzero $\beta_0$ and long times, because the universal term is small, and the branch constant of Appendix~\ref{app:eq3} must be pinned first.

    \item The boundary exponent from the first gap. The gap between the leading and first subleading eigenvalue measures the lightest boundary dimension,
    \begin{equation}
        \mathrm{Im}(\lambda_1-\lambda_0) = -\frac{\pi x_1}{v\,T},
        \label{eq:cft:eq4}
    \end{equation}
    the velocity converting the universal phase gap into $x_1$. Recovering $x_1\approx\tfrac12$ confirms the free boundary condition, that $\ket{X+}$ is the boundary state we think it is. More generally, every subleading eigenvalue is a boundary operator sitting at phase gap $\pi x/(vT)$: the spectrum is a ladder of \ac{CFT} states, and reading that ladder---including a part of it displaced by momentum $\pi$ once the model is interacting---is a central result of Section~\ref{sec:results}.

    \item Emergent dual unitarity. Equation~\eqref{eq:cft:transfer} also fixes how the spectrum evolves with $T$: the moduli of the leading eigenvalues approach a common value while their phases stay distinct,
    \begin{equation}
        \frac{|\mu_i(T)|}{|\mu_0(T)|}
            = 1 - \frac{a_i}{T} - \frac{b_i}{T^2} - \cdots
            \;\xrightarrow[T\to\infty]{}\; 1 .
        \label{eq:gapclose}
    \end{equation}
    A band of equal moduli with distinct phases is the fingerprint of a unitary operator up to an overall scale: the temporal transfer matrix becomes effectively unitary at long times. This is a consequence of conformal invariance, not of integrability, and it is the structure behind the logarithmic entropy growth above. It is also double-edged: the same degeneracy of moduli that signals the emergent conformal structure makes the leading eigenvectors progressively harder to resolve, which is what limits the reach of the method (Section~\ref{sec:numerics}). We verify it by tracing $\lambda_0(T)$ in the complex plane, where it draws a circle of near-constant radius---the non-universal modulus---with a winding phase.
\end{enumerate}

Section~\ref{sec:annni} first establishes, in equilibrium, that the ANNNI-type chain is critical and Ising-class, so that these predictions apply; the sections that follow build the tensor-network machinery to measure them through time.
